\newpage{\ }
\cleardoublepage

\thispagestyle{empty}
\chapter*{Agradecimientos}
\addcontentsline{toc}{chapter}{Agradecimientos} % Agregar a la TOC sin número

Esta tesis fue financiada por el Ministerio de Ciencia, Tecnología e Innovación Productiva (Argentina) y la Agencia Nacional de Promoción de la Investigación, el Desarrollo Tecnológico y la Innovación, mediante el subsidio PICT-2017-2142 otorgado a Thomas Kitzberger. 

Agradezco al Consejo Nacional de Investigaciones Científicas y Técnicas (CONICET, Argentina) por haberme otorgado una beca doctoral, a la Universidad Nacional del Comahue (UNCo) por permitirme estudiar esta carrera, y al Instituto de Investigaciones en Biodiversidad y Medioambiente (INIBIOMA, CONICET-UNCo) por brindarme el lugar de trabajo y la infraestructura necesaria.

Al Parque Nacional Nahuel Huapi (Administración de Parques Nacionales, Argentina) por permitirme realizar el estudio de campo correspondiente al capítulo 2 (permiso nro. 1680). 

A las personas e instituciones que desarrollan bases de datos y software de libre acceso. En particular, agradezco a los desarrolladores de Google Earth Engine, R, Stan, brms, mgcv, terra, y Rcpp. A todas las personas e instituciones que generan contenido educativo de libre acceso: autores de libros o artículos científicos, bloggers, youtubers, las generosas almas que responden preguntas en foros, y quienes se dedican a liberar contenido protegido por licencias. Gracias a su existencia pude aprender las herramientas necesarias para llevar a cabo esta tesis.  

A quienes participaron en los artículos que conforman los Capítulos 2 y 3 de esta tesis: Ana Cingolani, Juan Paritsis, Florencia Tiribelli, Mónica Mermoz y Luciana Ammassari, así como a las revisoras anónimas y al editor, Stacy Druri, que contribuyeron a mejorarlos. 

A Iván Renison, por revisar el Capítulo 4, por su ayuda con el paquete \texttt{FireSpread}, y por su paciencia para lidiar con mi código desprolijo para sus elevados estándares.

A Marcelo Bari, Sandro Lobos y Jorge Bonansea, quienes señalaron los puntos de ignición de muchos incendios, lo que permitió estimar el modelo de propagación. En particular, agradezco a Marcelo Bari por proveer la base de datos de focos confeccionada por el ICE del PNNH y por recordarme que el fuego no es sólo lo que vemos con satélites. A Fulden Batibeniz y Yann Quillcaille, quienes compartieron sus proyecciones diarias de FWI utilizadas en el Capítulo 4.

A Florencia Tiribelli, que en los comienzos de la tesis ofició de mentora y guía espiritual.

A Sofía Cingolani, quien colaboró con el trabajo de campo del Capítulo 2, revisó partes de la tesis, y tantas veces me ayudó a desenredar ideas.

A Ana Cingolani, mi codirectora latente, por su entusiasmo para discutir hasta el más mínimo detalle de las disquisiciones que le planteara. 

Al jurado, conformado por Sofía González, Ricardo Grau y Pablo Aceñolaza. Agradezco su lectura detenida; sus comentarios ayudaron a dejar el texto más claro, y aprecio que me hayan instado a resaltar aspectos relevantes de los resultados que había pasado por alto.

Y a mis directores, Thomas Kitzberger y Pájaro Morales (alias Juan Manuel), por brindarme la oportunidad de trabajar en este proyecto, propiciando un trato horizontal que fomenta la discusión. Y en particular, por su confianza; por darme la libertad de explorar mis intereses académicos a pesar de que no siempre se alinearan con los objetivos específicos de la tesis.